\documentclass[a4paper,11pt]{article}
\usepackage[utf8]{inputenc}
\usepackage[T1]{fontenc}
\usepackage[ngerman]{babel}
\usepackage{amsmath,amssymb}
\usepackage{xcolor}
\usepackage{colortbl}
\usepackage{array}
\usepackage{tikz}
\usepackage{enumitem}
\usepackage{geometry}
\geometry{a4paper,margin=2cm}
\definecolor{bgdark}{HTML}{1E1E1E}
\definecolor{fgtext}{HTML}{E6E6E6}
\definecolor{accent}{HTML}{6CB6FF}
\definecolor{accent2}{HTML}{F2A65A}
\definecolor{gridcol}{HTML}{4A4A4A}
\definecolor{linecol}{HTML}{8A8A8A}
\pagecolor{bgdark}
\color{fgtext}
\arrayrulecolor{linecol}
\newcommand{\karo}[1]{\par\noindent\begin{tikzpicture}\draw[step=5mm, gridcol, very thin] (0,0) grid (\linewidth,#1);\end{tikzpicture}\par}
\newcommand{\aufgabe}[2]{\medskip\par\noindent{\color{accent}\large\textbf{Aufgabe #1}}\hfill{\color{accent2}\small #2}\par\vspace{0.3em}\hrule height0.4pt\vspace{0.6em}}
\newcommand{\kopf}[1]{\noindent{\color{accent}\Large\textbf{Numerik — #1}}\par\vspace{0.4em}\noindent\textbf{Name:}~\rule[-2pt]{5cm}{0.4pt}\hfill\textbf{Datum:}~\rule[-2pt]{3cm}{0.4pt}\par\vspace{0.3em}\hrule height0.8pt\vspace{1em}}
\setlength{\parindent}{0pt}
\begin{document}
\kopf{Globale Optimierung}

\textbf{Notation der Vorlesung:} Eine 1D-Shekel-Funktion mit ``Foxholes'' lautet
\[
f(x)=-\sum_{i} \frac{c_i}{(x-a_i)^2+r_i},
\]
mit $a_i$ = Lage, $c_i$ = Tiefe und $r_i$ = Breite des $i$-ten Lochs. Lokale Optimierung sucht
$x^*=\arg\min_{x\in\mathcal N(x_0)} f(x)$ in einer Umgebung von $x_0$, globale Optimierung
$x^*=\arg\min_{x\in\Omega} f(x)$ auf dem ganzen Gebiet. Ist $f$ konvex, so ist jedes lokale
Minimum bereits global. Newton zur Optimierung: $x_{n+1}=x_n-\dfrac{f'(x_n)}{f''(x_n)}$,
modifiziert für Minima $x_{n+1}=x_n-\dfrac{f'(x_n)}{|f''(x_n)|}$. Zweite-Ableitung-Test:
$f''>0$ Minimum, $f''<0$ Maximum, $f''=0$ unklar.

\aufgabe{1}{leicht--mittel}
Gegeben sei die 1D-Shekel-Funktion mit zwei Foxholes
\[
f(x)=-\frac{2}{(x-0)^2+1}-\frac{4}{(x-5)^2+1}.
\]
(Loch A: $a=0,\,c=2,\,r=1$; Loch B: $a=5,\,c=4,\,r=1$.)
\begin{enumerate}[label=(\alph*)]
  \item Werten Sie $f$ an den Stellen $x=0$, $x=2{,}5$ und $x=5$ aus.
  \item Welche der Stellen $x=0$ und $x=5$ ist ein \emph{lokales}, welche das \emph{globale}
        Minimum? Begründen Sie über die Tiefe $c_i$.
\end{enumerate}
\karo{5cm}

\aufgabe{2}{mittel}
Newton zur Optimierung wird auf das Polynom
\[
f(x)=x^4-8x^2+3
\]
angewendet (zwei Minima, ein Maximum).
\begin{enumerate}[label=(\alph*)]
  \item Bilden Sie $f'(x)$ und $f''(x)$.
  \item Führen Sie ausgehend von $x_0=1{,}5$ \emph{einen} Newton-Schritt
        $x_1=x_0-\dfrac{f'(x_0)}{f''(x_0)}$ aus.
  \item Führen Sie einen \emph{zweiten} Schritt $x_2=x_1-\dfrac{f'(x_1)}{f''(x_1)}$ aus und
        vergleichen Sie mit dem exakten Minimum $x^*=2$.
\end{enumerate}
\karo{6cm}

\aufgabe{3}{leicht--mittel}
Betrachten Sie $f(x)=x^4-8x^2+3$ aus Aufgabe 2.
\begin{enumerate}[label=(\alph*)]
  \item Bestimmen Sie alle kritischen Punkte über $f'(x)=0$.
  \item Klassifizieren Sie jeden kritischen Punkt mit dem Zweite-Ableitung-Test
        ($f''>0$ Min, $f''<0$ Max, $f''=0$ unklar) als Minimum, Maximum oder Sattel.
\end{enumerate}
\karo{5cm}

\aufgabe{4}{mittel}
Der \emph{modifizierte} Newton-Schritt für Minimierung lautet
\[
x_{n+1}=x_n-\frac{f'(x_n)}{|f''(x_n)|}.
\]
\begin{enumerate}[label=(\alph*)]
  \item Zeigen Sie rechnerisch, dass die Schrittrichtung $d=x_{n+1}-x_n$ stets der Bedingung
        $f'(x_n)\cdot d\le 0$ genügt (also ``bergab'').
  \item Erklären Sie geometrisch, warum der Standard-Newton-Schritt
        $-f'/f''$ in einem Bereich mit $f''<0$ (konkav, Richtung Maximum) \emph{bergauf}
        laufen würde, der modifizierte Schritt aber nicht.
\end{enumerate}
\karo{6cm}

\aufgabe{5}{leicht--mittel}
Random Restarts: Bei jedem Neustart landet ein lokaler Optimierer mit Wahrscheinlichkeit $p$
im globalen Becken. Die Erfolgswahrscheinlichkeit nach $K$ unabhängigen Neustarts ist
$P=1-(1-p)^K$, die Umkehrung $K=\dfrac{\ln(1-P)}{\ln(1-p)}$.
\begin{enumerate}[label=(\alph*)]
  \item Berechnen Sie $P$ für $p=0{,}2$ und $K=10$.
  \item Wie viele Neustarts $K$ braucht man für $P=99\,\%$, wenn $p=0{,}05$?
  \item Vergleichen Sie: Wie verändert sich das nötige $K$, wenn $p$ von $0{,}2$ auf $0{,}05$ fällt?
\end{enumerate}
\karo{5cm}

\aufgabe{6}{mittel}
Lokale Methoden bei nicht-konvexen Funktionen.
\begin{enumerate}[label=(\alph*)]
  \item Erklären Sie, warum Newton bzw. Gradientenabstieg auf der Shekel aus Aufgabe 1 im
        flacheren lokalen Minimum bei $x=0$ ``stecken bleiben'' kann.
  \item Diskutieren Sie die Abhängigkeit vom Startpunkt $x_0$: Für welche Startwerte $x_0$
        konvergiert ein lokales Verfahren vermutlich zum globalen Minimum bei $x=5$, für welche
        zum lokalen bei $x=0$? (Sensitivität gegenüber Anfangsbedingungen.)
\end{enumerate}
\karo{6cm}

\aufgabe{7}{mittel--schwer}
Eine Shekel besitze drei Minima an den Stellen $x=0$, $x=4$ und $x=7$, wobei $x=4$ das globale
Minimum ist. Die trennenden lokalen Maxima liegen näherungsweise bei $x=2{,}1$ (zwischen $0$ und $4$)
und $x=5{,}6$ (zwischen $4$ und $7$). Das Suchgebiet ist $\Omega=[-1,\,9]$.
\begin{enumerate}[label=(\alph*)]
  \item Geben Sie für jedes der drei Minima das Einzugsgebiet (Basin of Attraction) als Intervall
        an (Grenzen = die trennenden Maxima bzw. die Ränder von $\Omega$).
  \item Bei einem uniformen Zufallsstart $x_0\sim\mathcal U([-1,9])$: Mit welcher Wahrscheinlichkeit
        landet ein lokaler Optimierer im globalen Becken (Minimum $x=4$)?
\end{enumerate}
\karo{6cm}

\aufgabe{8}{mittel}
Drei globale Strategien im Vergleich: \emph{Random Restarts}, \emph{Basin Hopping}
(lokale Optimierung + Zufallssprünge $\delta\sim\mathcal D$) und \emph{Simulated Annealing}
(Akzeptanz eines schlechteren Schritts mit Wahrscheinlichkeit $\exp(-\Delta f/T)$, Temperatur $T$ sinkt).
\begin{enumerate}[label=(\alph*)]
  \item Beschreiben Sie für jedes Verfahren in einem Satz, \emph{wie} es lokalen Minima entkommt.
  \item Berechnen Sie die Simulated-Annealing-Akzeptanzwahrscheinlichkeit $\exp(-\Delta f/T)$
        für eine Verschlechterung $\Delta f=2$ bei den Temperaturen $T=5$ und $T=0{,}5$.
  \item Was bedeutet das Ergebnis aus (b) für das ``Abkühlen'' ($T\to 0$)?
\end{enumerate}
\karo{6cm}

\aufgabe{9}{schwer}
Konvex vs. nicht-konvex. Eine Funktion heißt konvex, falls für alle $x,y$ und $\lambda\in[0,1]$
\[
f(\lambda x+(1-\lambda)y)\le \lambda f(x)+(1-\lambda)f(y).
\]
\begin{enumerate}[label=(\alph*)]
  \item Prüfen Sie über $f''$, ob $f(x)=x^2+3$ konvex ist, und folgern Sie, ob ein lokales
        Verfahren hier garantiert das globale Minimum findet.
  \item Prüfen Sie $g(x)=x^4-8x^2+3$ (aus Aufgabe 2) auf Konvexität, indem Sie $g''(x)$ auf
        Vorzeichenwechsel untersuchen. Was folgt daraus für die Eignung lokaler Methoden?
  \item Weisen Sie für $h(x)=x^4-8x^2+3$ die Verletzung der Konvexitäts-Ungleichung an einem
        konkreten Paar nach, z.\,B. $x=-2$, $y=2$, $\lambda=\tfrac12$.
\end{enumerate}
\karo{7cm}

\end{document}
